\section{Analysis: from scenario to measured result}\label{sec:analysis}
Table~\ref{tab:thread} closes the loop from each scenario requirement (Table~\ref{tab:req})
through the architectural mechanism that serves it to the evidence that it works.

\begin{table}[H]\centering\small
\caption{Traceability: requirement $\rightarrow$ design $\rightarrow$ traffic $\rightarrow$ measured result.}\label{tab:thread}
\begin{tabularx}{\linewidth}{>{\raggedright}p{3.0cm}>{\raggedright}p{3.6cm}>{\raggedright}p{2.5cm}>{\raggedright\arraybackslash}X}
\toprule
Requirement & Mechanism (\S\ref{sec:slices}, \S\ref{sec:arch}) & Traffic & Evidence\\
\midrule
Three services must be logically separated & 3 S-NSSAIs, NSSF/AMF/SMF slice support, per-UE subscription & -- & Allowed/Rejected NSSAI per UE (Tab.~\ref{tab:nssai}); unsubscribed slice refused (E1)\\
Each slice has its own anchor and addressing & DNN$\rightarrow$UPF selection, per-DNN pools and tun & ICMP per session & GTP-U to 10.10.0.1 vs 10.10.0.3 (Lst.~\ref{lst:gtpu}); PFCP 5$\times$UPF-A, 2$\times$UPF-B\\
UE reaches the external DN & N3/N6 namespaces, DN routes, NAT & ping, iperf3, HTTP & 7/7 sessions, 0\,\% loss; 138/18.7/2.1\,Mbit/s; HTTP 200\\
eMBB high rate but capped & HTB MBR + Session-AMBR & on/off TCP & goodput at 92--94\,\% of cap (E1, E3, E4)\\
URLLC p99\,$<$\,10\,ms and lossless under eMBB load & dedicated UPF-B, MBR on eMBB & 100\,Hz probe & p99 0.86\,ms, worst 1.43\,ms, 0\,\% loss while eMBB at cap (E3)\\
mIoT reliable delivery & shared UPF-A, low AMBR & Poisson, 400--1000 devices & 0\,\% loss, ACK p99 $\approx$1--2\,ms (E3)\\
Operator can retune a slice live & tc class/qdisc change via sandbox & all slices & MBR and delay changes take effect within 1\,s, other slices unaffected (E4)\\
Dedicated URLLC user plane worthwhile & UPF-B vs.\ shared UPF-A & UDP overload & median p99 7.9 vs.\ 12.9\,ms; 72\,\% vs.\ 35\,\% of seconds within target (E5)\\
\bottomrule
\end{tabularx}
\end{table}

\paragraph{What the results mean for the scenario.} Within its normal operating envelope
(eMBB at or below its MBR), the factory network meets every requirement of
Table~\ref{tab:req}. The decisive mechanism is the \emph{rate cap on the aggressive slice}.
With eMBB held at its MBR, URLLC and mIoT saw no loss and sub-2\,ms p99 even when eMBB was
saturating (E3). Once that cap is removed and the platform is overloaded (E5), every slice
suffers, because the gNB, the RLS link and the CPU are shared. A dedicated UPF then still
shifts the URLLC latency distribution down, but it does not prevent loss. The practical
lesson is that slice isolation must be enforced at \emph{every} shared resource. In this
testbed those are the transport (tc) and the UPF. In a real network the RAN scheduler is the
most important one, and it is precisely what a software-only RAN cannot model
(\S\ref{sec:limits}).

\paragraph{Signalling cost.} One registration plus one PDU session costs about 12 NGAP
messages, 4 PFCP messages and about 25 distinct SBI requests per UE. The capture holds 82
NGAP frames and 354 non-NRF SBI requests for 7 UEs, the latter seen twice via the SCP. That
amounts to 20--30\,ms of core processing on this platform (Table~\ref{tab:sigtime}). For a
multi-slice UE the PDU procedure runs once per S-NSSAI. Both requests are handled in
parallel: the two PFCP sessions on UPF-A and UPF-B are created within 0.6\,ms of each other
(Fig.~\ref{fig:ladder}). The second N2 response arrives about 200\,ms later only because of
UERANSIM's internal timer.

\paragraph{Reading the latency numbers.} The absolute RTTs (0.3--1.1\,ms median) are the
core and transport cost only; they include no air interface. They are also modulated by the
VM's CPU state (E7): an idle VM looks \emph{slower} than a busy one, so we only compare
measurements taken under like-for-like load. For URLLC dimensioning, the relevant results
are the relative ones: the 10\,ms increase produced by a 5\,ms/direction transport impairment
(E4), and the distribution shift caused by UPF sharing (E5).
